\paragraph{¿Por qué modelos basados en agentes?}

El modelado basado en agentes (ABM, \emph{Agent-Based Modeling}) ofrece ventajas estructurales frente a los enfoques alternativos para simular la propagación de narrativas en redes sociales.

Los modelos de ecuaciones diferenciales ordinarias (ODE), heredados de la epidemiología (modelos SIR, SIS), trabajan con densidades poblacionales homogéneas y asumen que todos los individuos son equivalentes y están en contacto potencial con todos los demás. Esta hipótesis de mezcla homogénea es fundamentalmente incorrecta para redes sociales, donde la estructura de la red ---hubs, comunidades, puentes débiles--- es precisamente el mecanismo que gobierna la difusión \cite{Watts2002,CentolaMacy2007}. Los ABM no imponen dicha homogeneidad: la heterogeneidad emerge de la interacción local entre agentes sobre una topología específica.

Los modelos estadísticos y de aprendizaje automático son potentes para predecir patrones observados, pero no son causales: describen correlaciones entre variables pasadas sin modelar el mecanismo generador. Para el objetivo de este trabajo ---medir cómo la topología y la psicología individual determinan la resiliencia cognitiva--- se requiere un modelo \emph{generativo} que permita intervenir en parámetros individuales y observar el efecto sobre la dinámica colectiva.

Los juegos evolutivos y los modelos de dinámica de opinión (DeGroot, Hegselmann-Krause) capturan la actualización de creencias pero tipicamente asumen homogeneidad de la red, ausencia de contenido emocional en los mensajes, o reglas de actualización simétricas incompatibles con la dirección asimétrica de las plataformas digitales.

El ABM resuelve estas limitaciones con tres propiedades: (i) permite asignar a cada agente un umbral $\theta_i$ individual muestreado de una distribución heterogénea, condición necesaria para observar cascadas no lineales \cite{Granovetter1978}; (ii) sitúa a cada agente en un nodo del grafo y hace que la difusión dependa explícitamente de la estructura de vecindad; y (iii) permite registrar trazas causales de cada mensaje, requisito imprescindible para el análisis \emph{post hoc} de la dinámica de resistencia.

El coste del ABM es computacional: cada réplica simula $N = 5\,000$ agentes durante 45 pasos de simulación, generando decenas de miles de mensajes por corrida. El diseño OFAT con 25 réplicas Monte Carlo y 700 runs totales requiere aproximadamente 30 minutos de CPU con implementación eficiente en Python/Mesa 3.x.

\medskip

El simulador implementa tres modelos de agente de complejidad psicológica creciente.
Su diseño responde a una lógica de control experimental: el modelo más simple actúa como
línea base estadística, el modelo intermedio concentra la potencia explicativa de los
experimentos y el modelo más rico constituye una propuesta teórica para trabajo futuro.
Los tres comparten la misma interfaz de paso de mensaje sobre el grafo, lo que permite
intercambiarlos sin modificar la lógica de la simulación.

\paragraph{Agente estocástico (control nulo).}
La clase \texttt{StochasticAgent} es un agente \emph{reactivo puro}: carece de estado
cognitivo interno y no mantiene ninguna representación de las narrativas que circulan por
la red. Su comportamiento en cada paso de la simulación se reduce a dos acciones
independientes parametrizadas por sendas probabilidades: con probabilidad $p_{\text{fwd}}$
reenvía el último mensaje recibido a un subconjunto aleatorio de sus seguidores en el
grafo, y con probabilidad $p_{\text{fire}}$ genera un mensaje de ruido neutro y lo emite
igualmente a un subconjunto aleatorio. Ambas acciones son independientes en cada paso, de
modo que un agente puede reenviar y generar ruido simultáneamente, solo una de las dos, o
ninguna.

La utilidad de este modelo no es descriptiva sino metodológica. Al eliminar toda
psicología individual, aisla la contribución de la \emph{topología de la red} y del
\emph{azar} en la difusión de contenidos. Cualquier diferencia observada entre las métricas
de este agente y las de los modelos más ricos debe atribuirse, por diferencia, a los
mecanismos cognitivos añadidos en estos últimos. Sin este control nulo sería imposible
separar el efecto estructural del grafo del efecto de los procesos individuales de
resistencia o persuasión.

\paragraph{Agente resistente (umbral lineal cognitivo).}
La clase \texttt{ResistantAgent} es el modelo implementado en los experimentos de esta
memoria. Incorpora un estado interno compuesto por dos escalares: la \emph{convicción}
$c_i \in [0,1]$, que cuantifica en qué medida el agente $i$ ha interiorizado una narrativa
concreta, y el \emph{umbral de resistencia} $\theta_i \in [0,1]$, que marca el nivel de
convicción necesario para que el agente adopte dicha narrativa. El umbral $\theta_i$ se
muestrea de una distribución bimodal que modela dos subpoblaciones —susceptible y
resistente—; su expresión formal se detalla en el diseño de la solución
(ecuación~\ref{eq:theta}).

Un mensaje pertenece a la narrativa rastreada si y solo si su veracidad $v$ satisface
$v \leq v_{\text{thr}}$, donde $v_{\text{thr}}$ es un umbral de veracidad global
configurable. Cuando el agente $i$ recibe un mensaje de la narrativa procedente del
agente $j$, actualiza su convicción según la regla:

\begin{equation}
  c_i \;\leftarrow\; c_i + (1 - c_i)\cdot\tau_{j \to i}\cdot\lambda_m\cdot\gamma,
  \label{eq:conviction}
\end{equation}

donde $\tau_{j \to i}$ es la confianza que el agente $i$ deposita en el emisor $j$ en la
capa del mensaje, $\lambda_m \in [0,1]$ es la \emph{carga emocional} del mensaje y
$\gamma > 0$ es la ganancia de persuasión por exposición. La forma multiplicativa de la
ecuación garantiza que $c_i$ permanezca acotada en $[0,1]$ y que la tasa de incremento
decrezca conforme la convicción se aproxima a la saturación, en línea con el marco del
modelo de umbral lineal de~\cite{KempeKleinbergTardos2003}. El agente adopta la narrativa
en el instante en que $c_i \geq \theta_i$.

La heterogeneidad de los umbrales $\{\theta_i\}$ —no su valor medio— es el factor
determinante en la formación de cascadas de adopción: una distribución concentrada impide
o universaliza la propagación, mientras que la bimodalidad introduce puntos de inflexión
no lineales en la dinámica de difusión~\cite{Granovetter1978,Watts2002}. Adicionalmente,
el hecho de que la adopción requiera exposición reiterada al mismo contenido convierte
este proceso en un \emph{contagio complejo}, cualitativamente distinto de la difusión
simple y más sensible a la estructura de los puentes débiles~\cite{CentolaMacy2007}.

Una vez adoptada la narrativa, el agente no se vuelve pasivo: en cada paso ulterior
re-emite espontáneamente su última narrativa con probabilidad $f$ (período infeccioso
superior a un único paso) y emite ruido neutro con probabilidad $p_{\text{fire}}$. Esta
capacidad de re-emisión autónoma genera un \emph{contagio autosostenido} en el que los
nodos adoptados actúan como fuentes secundarias, amplificando la dinámica de cascada más
allá de la fuente original.

\paragraph{Agente cognitivo-emocional (bayesiano).}
Las clases \texttt{BayesianAgent} y \texttt{EmotionalBrain} forman el modelo teórico más
rico del simulador. No se incluye en los experimentos de esta memoria, pero su diseño
ilustra una dirección de extensión naturalista bien fundamentada. El agente ejecuta el
bucle \emph{OODA} (observar, orientar, decidir, actuar)~\cite{RushingHerschXu2026}, en el
que la fase de orientación está mediada por un \emph{cerebro emocional} vectorial.

El estado interno del cerebro es un vector $\mathbf{e} \in [0,1]^{11}$: las ocho emociones
básicas del modelo de Plutchik~\cite{Plutchik1980} —alegría, confianza, miedo, sorpresa,
tristeza, aversión, ira y anticipación— más tres dimensiones metacognitivas: atención,
credulidad y autoconfianza. Cada agente posee además un vector de \emph{genética}
$\mathbf{g} \in [0,1]^{11}$ que actúa como amplificador estable e individual de la
respuesta emocional, modelando la variabilidad interindividual en sensibilidad afectiva.

Al observar un mensaje, el cerebro actualiza su estado emocional siguiendo el principio
de codificación predictiva y actualización bayesiana de creencias~\cite{Friston2010}:

\begin{equation}
  \Delta = \mathbf{m}_{\text{vec}} \cdot g_{\text{mod}} \cdot \mathbf{g}
           \cdot \mathrm{attn} \cdot \tau \cdot \mathrm{cred} \cdot \mathrm{circ}(t),
  \qquad
  \mathbf{e} \;\leftarrow\; \operatorname{clip}(\mathbf{e} + \Delta,\; 0,\; 1),
  \label{eq:brain}
\end{equation}

donde $\mathbf{m}_{\text{vec}}$ es el vector emocional del mensaje, $g_{\text{mod}}$ un
factor de ganancia por modalidad (texto, imagen, vídeo), $\mathrm{attn}$ la atención
actual del agente, $\tau$ la confianza en el emisor, $\mathrm{cred}$ la credulidad
disposicional y $\mathrm{circ}(t)$ un modulador circadiano. Entre pasos se aplica un
decaimiento exponencial por dimensión
$e_k \leftarrow e_k \cdot 2^{-1/h_{1/2}}$,
donde $h_{1/2}$ es la vida media emocional en pasos de simulación. El modulador circadiano
adopta la forma
\[
  \mathrm{circ}(t) = \mathrm{base} + \mathrm{amp}
  \cdot \cos\!\Bigl(\frac{2\pi(h - h_{\mathrm{pico}})}{24}\Bigr),
\]
con un pico de sugestionabilidad positiva en horas matinales y un pico de afecto negativo
nocturno, en consonancia con los patrones afectivos observados en redes sociales a lo
largo del día~\cite{GolderMacy2011}.

El motivo por el que este modelo no entra en los experimentos es de orden metodológico: el
determinismo y la trazabilidad del agente de umbral bastan para medir la resistencia a la
desinformación de forma reproducible y controlada. El cerebro emocional vectorial introduce
fuentes de no-determinismo adicionales —el decaimiento por dimensión, la modulación
circadiana continua, la interacción multiplicativa de once dimensiones— que, sin un diseño
experimental específicamente orientado a ellas, oscurecerían la atribución causal y
elevarían el coste computacional sin una ganancia proporcional en la validez interna del
experimento.

\begin{figure}[h!]
\centering
\fbox{\parbox{0.85\linewidth}{\centering\vspace{1.5cm}
PENDIENTE (figura a generar): Esquema del cerebro emocional de 11~dimensiones del
\texttt{EmotionalBrain}. Las 8~emociones de Plutchik aparecen en la periferia y las
3~dimensiones metacognitivas (atención, credulidad, autoconfianza) en el centro. Flechas
indican el flujo de actualización al recibir un mensaje: entrada $\mathbf{m}_{\text{vec}}$,
modulación por $\mathbf{g}$, $\tau$, $\mathrm{circ}(t)$, y salida $\mathbf{e}$
actualizado mediante la ecuación~\eqref{eq:brain}.%
\vspace{1.5cm}}}
\caption{Arquitectura del cerebro emocional vectorial del agente cognitivo-emocional.
         El estado interno es un vector de 11~dimensiones actualizado en cada observación
         (ecuación~\eqref{eq:brain}).}
\label{fig:cerebro}
\end{figure}

\begin{figure}[h!]
\centering
\fbox{\parbox{0.85\linewidth}{\centering\vspace{1.5cm}
PENDIENTE (figura a generar): Rueda de emociones de Plutchik~\cite{Plutchik1980} con las
ocho emociones básicas y sus combinaciones diádicas primarias, coloreada según valencia
afectiva (tonos cálidos para emociones positivas, fríos para negativas). Se destacan las
ocho dimensiones adoptadas en el modelo:
\textit{joy, trust, fear, surprise, sadness, disgust, anger, anticipation}.%
\vspace{1.5cm}}}
\caption{Rueda de emociones de Plutchik~\cite{Plutchik1980}. Las ocho emociones básicas
         de la periferia constituyen las primeras ocho dimensiones del vector de estado
         $\mathbf{e}$ del \texttt{EmotionalBrain}.}
\label{fig:plutchik}
\end{figure}

\begin{table}[h!]
\centering
\small
\caption{Comparativa de los tres modelos de agente implementados en el simulador.}
\label{tab:agentes}
\begin{tabular}{@{}p{3cm} p{2.7cm} p{4.1cm} p{1.6cm} p{2cm}@{}}
\toprule
\textbf{Modelo} &
\textbf{Estado interno} &
\textbf{Regla de decisión} &
\textbf{Determinista} &
\textbf{En experimentos} \\
\midrule
Estocástico
(\texttt{StochasticAgent})
  & Ninguno
  & Reenvío y ruido con probabilidades $p_{\text{fwd}}$, $p_{\text{fire}}$
  & No
  & Sí (control) \\
\addlinespace
Resistente
(\texttt{ResistantAgent})
  & $c_i,\,\theta_i \in [0,1]$
  & Umbral lineal: adoptar si $c_i \geq \theta_i$ \cite{KempeKleinbergTardos2003,Granovetter1978}
  & Sí
  & Sí \\
\addlinespace
Cognitivo-emocional
(\texttt{BayesianAgent})
  & $\mathbf{e},\,\mathbf{g} \in [0,1]^{11}$
  & Actualización bayesiana del vector emocional \cite{Friston2010,Plutchik1980}
  & No
  & No \\
\bottomrule
\end{tabular}
\end{table}
